\section{Легенда, или Постановка задачи}

Летним вечером двое мальчиков сидели на лавочке, смотрели на звёзды и спорили,
сколько атомов поместится во Вселенной. Они называли всё большие числа --- миллион,
миллиард, триллион --- и дошли до \emph{ундециллиона}, до \(10^{36}\). Им
показалось мало.

Тогда они поступили так, как поступает всякий математик, столкнувшись с числом, для
которого ещё нет имени: они дали имя. Первый шаг они сделали сразу большим: пусть
\(a = \avalue\) --- ундециллион, возведённый в самого себя. А дальше
--- по восходящей: каждое следующее число есть предыдущее, возведённое в самого себя.
Так родился ряд, ведущий к \textbf{крипто коме}.

\begin{figure}[h]
\centering
\begin{tikzpicture}[
  node distance=8mm and 14mm,
  lvl/.style={draw,rounded corners,minimum width=13mm,minimum height=9mm,font=\small},
  arr/.style={-{Stealth[length=2mm]},thick},
]
  \node[lvl] (a) {$a$};
  \node[lvl,right=of a] (b) {$b$};
  \node[lvl,right=of b] (c) {$c$};
  \node[right=of c] (dots) {$\cdots$};
  \node[lvl,right=of dots] (z) {$z$};
  \node[lvl,fill=black!8,right=of z] (cc) {$\ccoma$};
  \draw[arr] (a) -- node[above,font=\scriptsize]{$(\cdot)^{(\cdot)}$} (b);
  \draw[arr] (b) -- (c);
  \draw[arr] (c) -- (dots);
  \draw[arr] (dots) -- (z);
  \draw[arr] (z) -- node[above,font=\scriptsize]{$z^{z}$} (cc);
\end{tikzpicture}
\caption{Лестница уровней: \(a\) --- первый шаг, \(z\) --- двадцать шестой,
крипто кома \(\ccoma = z^{z}\) --- её вершина. Всего уровней: \nlevels.}
\label{fig:ladder}
\end{figure}
